\section{为什么本章必须避免伪严谨}

手稿需要一种形式语言，因为后续章节要讨论状态形成、状态转变与干预设计。但这种语言必须诚实地说明它属于哪一种形式化。目前，这个项目\emph{并没有}一个已经拟合好的定量模型，也没有为人类日常行为实验性识别出系数。因此，它不应当把随意写出的公式包装成已经成立的数学或已定论的实证规律。

本章采取了一个更严格的立场。它提供的是：
\begin{enumerate}[nosep]
  \item 一套用于描述行为状态的有纪律词汇；
  \item 一个用于比较阻力来源的启发式评分支架；
  \item 关于障碍与准备的定性单调假设；
  \item 对这套记号能主张什么、不能主张什么所设定的明确证据边界。
\end{enumerate}

\begin{discussion}
本章的来源支持主要来自两个方向。\textbf{教材基础层}来自 Eugene Izhikevich 的 \emph{Dynamical Systems in Neuroscience}、Michael Brin 与 Garrett Stuck 的 \emph{Introduction to Dynamical Systems}，以及 Steven Strogatz 的 \emph{Nonlinear Dynamics and Chaos} 等动力系统教材。这些来源为使用状态空间、稳定性与转变语言作为数学脚手架提供了正当性。\textbf{生物学正当性}则来自后续实证章节及其参考来源，它们表明结构化的状态转变与稳定性并不只是比喻。它们\emph{没有}提供的，是一个能够刻画日常行动失败的、已校准的标量定律。
\end{discussion}